\documentclass[12pt,a4paper]{article}

% Package imports
\usepackage[utf8]{inputenc}
\usepackage[margin=1in]{geometry}
\usepackage{graphicx}
\usepackage{amsmath}
\usepackage{amssymb}
\usepackage{enumitem}
\usepackage{xcolor}
\usepackage{listings}
\usepackage{tcolorbox}
\usepackage{titlesec}
\usepackage{hyperref}
\usepackage{fancyhdr}

% Header and footer
\pagestyle{fancy}
\fancyhf{}
\rhead{World Balance Architects}
\lhead{Game Design Document}
\rfoot{Page \thepage}

% Color definitions
\definecolor{codebackground}{rgb}{0.95,0.95,0.95}
\definecolor{sectioncolor}{rgb}{0.1,0.3,0.5}

% Section formatting
\titleformat{\section}
  {\normalfont\LARGE\bfseries\color{sectioncolor}}
  {\thesection}{1em}{}

\titleformat{\subsection}
  {\normalfont\Large\bfseries\color{sectioncolor}}
  {\thesubsection}{1em}{}

% Code listing style
\lstset{
    backgroundcolor=\color{codebackground},
    basicstyle=\ttfamily\small,
    breaklines=true,
    frame=single,
    tabsize=2
}

% Title information
\title{
    \Huge\textbf{World Balance Architects} \\
    \vspace{0.5cm}
    \Large AI-Powered Environmental Strategy Game \\
    \vspace{0.3cm}
    \large Complete Design Document
}
\author{Game Design Specification}
\date{January 2026}

\begin{document}

\maketitle
\thispagestyle{empty}

\newpage
\tableofcontents
\newpage

% ================================================================
% SECTION 1: CORE GAME DESIGN
% ================================================================
\section{The Simplest Version of the Combined Game}

This section outlines the fundamental architecture of the game, designed specifically to be suitable for AI laboratory implementation while maintaining engaging gameplay mechanics.

\subsection{World Structure (Grid-Based System)}

The game world is represented as a two-dimensional grid-based map with recommended dimensions of 25 cells wide by 18 cells high. This grid-based approach provides discrete decision spaces that are ideal for AI algorithm implementation and comparison.

\subsubsection{Cell Types}

Each cell in the grid can be assigned one of the following terrain types:

\begin{itemize}[leftmargin=2cm]
    \item \textbf{Land} -- Basic terrain type; can be converted to other types through player actions
    \item \textbf{River} -- Water source that flows through the map; critical for ecosystem balance
    \item \textbf{Farm} -- Agricultural land used for food production; requires water resources
    \item \textbf{Forest} -- Natural ecosystem that produces oxygen and regulates temperature
    \item \textbf{Reservoir} -- Water storage facility that captures and distributes water to nearby areas
    \item \textbf{City/Energy} -- Optional infrastructure for energy production and resource management
\end{itemize}

\subsection{Global Planet Metrics}

The game features four primary planetary health indicators that are constantly monitored and displayed on the user interface. These meters represent the overall health and sustainability of the planet:

\subsubsection{Water Level (0--100)}

Represents the total available water resources on the planet. This metric is influenced by:
\begin{itemize}
    \item River flow and distribution
    \item Reservoir storage capacity
    \item Farm irrigation consumption
    \item Natural water cycle simulation
\end{itemize}

\subsubsection{Food Supply (0--100)}

Indicates the amount of food available to sustain the planet's population. Key factors include:
\begin{itemize}
    \item Active farm production
    \item Successful crop harvests
    \item Water availability for agriculture
    \item Agricultural infrastructure development
\end{itemize}

\subsubsection{Oxygen Production (0--100)}

Measures the atmospheric oxygen content produced by natural ecosystems. This is affected by:
\begin{itemize}
    \item Forest coverage and density
    \item Balance between vegetation and deforestation
    \item Overall ecosystem health
\end{itemize}

\subsubsection{Temperature Regulation (0--100)}

Tracks the planetary temperature with the goal of maintaining values in the middle range (ideally 40--60). Temperature is influenced by:
\begin{itemize}
    \item Industrial activity and energy production
    \item Forest coverage (cooling effect)
    \item Water distribution across the map
    \item Agricultural development intensity
\end{itemize}

\subsection{Turn-Based Game Loop}

The game operates on a turn-based system, which is particularly well-suited for AI laboratory implementation because it allows for discrete, comparable decision-making processes.

\subsubsection{Turn Sequence}

Each complete turn follows this precise sequence:

\begin{enumerate}[leftmargin=2cm]
    \item \textbf{Agent A Decision Phase}
    \begin{itemize}
        \item Agent A analyzes the current world state
        \item Selects one action from available options
        \item Action is applied to the world grid
    \end{itemize}
    
    \item \textbf{World Update Phase (After Agent A)}
    \begin{itemize}
        \item Water flow simulation executes
        \item Crop growth is calculated
        \item Oxygen production is updated
        \item Temperature changes are computed
        \item All global meters are recalculated
    \end{itemize}
    
    \item \textbf{Agent B Decision Phase}
    \begin{itemize}
        \item Agent B analyzes the current world state (including Agent A's changes)
        \item Selects one action from available options
        \item Action is applied to the world grid
    \end{itemize}
    
    \item \textbf{World Update Phase (After Agent B)}
    \begin{itemize}
        \item Water flow simulation executes
        \item Crop growth is calculated
        \item Oxygen production is updated
        \item Temperature changes are computed
        \item All global meters are recalculated
    \end{itemize}
    
    \item \textbf{Turn Completion}
    \begin{itemize}
        \item Turn counter increments
        \item System checks for win/loss conditions
        \item Next turn begins
    \end{itemize}
\end{enumerate}

\subsubsection{Advantages for AI Laboratory Research}

This turn-based approach provides several critical advantages for AI research:

\begin{itemize}
    \item \textbf{Discrete Decision Space} -- Each action can be clearly defined and analyzed
    \item \textbf{Reproducible Results} -- Same starting conditions always produce same outcomes
    \item \textbf{Algorithm Comparison} -- Different AI approaches can be directly compared
    \item \textbf{State Analysis} -- Game state can be frozen and examined at any point
    \item \textbf{Performance Measurement} -- Clear metrics for evaluating AI effectiveness
\end{itemize}

% ================================================================
% SECTION 2: ACTIONS
% ================================================================
\section{Actions (Keeping It Interesting, Not Overwhelming)}

Each AI agent is provided with a carefully curated set of actions designed to create strategic depth without overwhelming complexity. The recommended action set consists of 5--7 distinct actions across multiple categories.

\subsection{Design Philosophy}

The action set is intentionally limited to:
\begin{itemize}
    \item Ensure each decision is meaningful and strategic
    \item Prevent analysis paralysis for both AI and human players
    \item Create clear trade-offs between different strategies
    \item Maintain computational feasibility for AI algorithms
\end{itemize}

\subsection{Action Categories and Descriptions}

\subsubsection{River and Water Flow Actions}

These actions manipulate water distribution across the planet:

\begin{tcolorbox}[colback=blue!5,colframe=blue!50!black,title=Build Canal]
\textbf{Description:} Converts a land cell into a river cell, creating new water pathways.

\textbf{Strategic Use:}
\begin{itemize}
    \item Redirects water flow to specific areas
    \item Supplies water to distant farms
    \item Creates water networks for ecosystem balance
\end{itemize}

\textbf{Cost Considerations:} Requires eco points and may have construction time
\end{tcolorbox}

\begin{tcolorbox}[colback=blue!5,colframe=blue!50!black,title=Place Reservoir]
\textbf{Description:} Constructs a water storage facility that captures and distributes water to nearby areas.

\textbf{Strategic Use:}
\begin{itemize}
    \item Stores water for drought periods
    \item Provides consistent irrigation to adjacent farms
    \item Stabilizes water availability in specific regions
    \item Acts as a buffer against water flow fluctuations
\end{itemize}

\textbf{Cost Considerations:} Higher initial cost but provides long-term benefits
\end{tcolorbox}

\begin{tcolorbox}[colback=blue!5,colframe=blue!50!black,title=Open/Close Gate]
\textbf{Description:} Toggles a gate tile that controls water flow through specific pathways.

\textbf{Strategic Use:}
\begin{itemize}
    \item Dynamically controls water distribution
    \item Prevents flooding in certain areas
    \item Redirects water during different seasons or conditions
    \item Allows for tactical water management
\end{itemize}

\textbf{Cost Considerations:} Low cost action that provides strategic flexibility
\end{tcolorbox}

\subsubsection{Ecosystem Management Actions}

These actions directly affect the planet's natural ecosystems:

\begin{tcolorbox}[colback=green!5,colframe=green!50!black,title=Plant Forest]
\textbf{Description:} Converts a land cell into a forest, creating a natural oxygen producer and temperature regulator.

\textbf{Strategic Use:}
\begin{itemize}
    \item Increases oxygen production over time
    \item Reduces planetary temperature
    \item Creates natural carbon sinks
    \item Improves overall ecosystem stability
\end{itemize}

\textbf{Long-term Benefits:}
\begin{itemize}
    \item Forests mature over several turns, becoming more effective
    \item Provides sustained environmental benefits
    \item Acts as a buffer against temperature extremes
\end{itemize}

\textbf{Cost Considerations:} Medium cost with delayed but significant returns
\end{tcolorbox}

\begin{tcolorbox}[colback=red!5,colframe=red!50!black,title=Clear Forest]
\textbf{Description:} Removes a forest cell, converting it back to land.

\textbf{Strategic Use:}
\begin{itemize}
    \item Clears space for farms or infrastructure
    \item Provides immediate resources or points
    \item Can be necessary for strategic repositioning
\end{itemize}

\textbf{Negative Consequences:}
\begin{itemize}
    \item Reduces oxygen production
    \item Increases planetary temperature
    \item Decreases ecosystem stability
    \item May trigger environmental penalties
\end{itemize}

\textbf{Cost Considerations:} Low immediate cost but high environmental cost
\end{tcolorbox}

\subsubsection{Agricultural Actions}

These actions focus on food production:

\begin{tcolorbox}[colback=yellow!5,colframe=yellow!50!black,title=Plant Crop]
\textbf{Description:} Establishes a farm on a land cell to begin food production.

\textbf{Strategic Use:}
\begin{itemize}
    \item Initiates food production cycle
    \item Must be placed near water sources
    \item Requires ongoing water supply to grow
\end{itemize}

\textbf{Growth Requirements:}
\begin{itemize}
    \item Adjacent or nearby water source
    \item Several turns to reach harvestable state
    \item Consistent water availability
\end{itemize}

\textbf{Cost Considerations:} Medium cost with turn-delayed returns
\end{tcolorbox}

\begin{tcolorbox}[colback=yellow!5,colframe=yellow!50!black,title=Harvest Crop]
\textbf{Description:} Collects food from a mature farm, adding to the food supply.

\textbf{Strategic Use:}
\begin{itemize}
    \item Converts matured crops into food points
    \item Increases global food meter
    \item Provides immediate scoring benefits
\end{itemize}

\textbf{Resource Impact:}
\begin{itemize}
    \item Consumes water resources
    \item May temporarily reduce farm productivity
    \item Requires farm to regrow before next harvest
\end{itemize}

\textbf{Cost Considerations:} Low action cost but consumes water resources
\end{tcolorbox}

\subsubsection{Energy Production Actions (Optional but Engaging)}

These actions provide energy while affecting other planetary systems:

\begin{tcolorbox}[colback=orange!5,colframe=orange!50!black,title=Build Solar Plant]
\textbf{Description:} Constructs a solar energy facility that generates power for the planet.

\textbf{Strategic Use:}
\begin{itemize}
    \item Provides renewable energy production
    \item Enables advanced infrastructure
    \item Supports technological development
\end{itemize}

\textbf{Environmental Impact:}
\begin{itemize}
    \item Slightly increases local temperature (heat absorption)
    \item Occupies land that could be used for farms or forests
    \item May reflect light affecting nearby ecosystems
\end{itemize}

\textbf{Cost Considerations:} High initial cost with sustained energy returns
\end{tcolorbox}

\subsection{Action Cost System}

To prevent action spamming and encourage strategic planning, each action requires "eco points" or equivalent resource cost:

\begin{itemize}
    \item \textbf{Low Cost Actions} (1--2 points): Gate toggle, harvest crop
    \item \textbf{Medium Cost Actions} (3--5 points): Plant forest, plant crop, build canal
    \item \textbf{High Cost Actions} (6--10 points): Place reservoir, build solar plant, clear forest (environmental penalty)
\end{itemize}

Players earn eco points through:
\begin{itemize}
    \item Maintaining balanced planetary metrics
    \item Achieving sustainability milestones
    \item Turn-based passive generation
    \item Completing environmental objectives
\end{itemize}

% ================================================================
% SECTION 3: SIMULATION RULES
% ================================================================
\section{Simulation Rules (Simple But Strategically Deep)}

The simulation engine does not require real-world physics precision. Instead, it employs simple, intuitive rules that create emergent complexity and strategic depth.

\subsection{Design Philosophy}

The simulation rules are designed to:
\begin{itemize}
    \item Be computationally efficient for AI planning
    \item Create clear cause-and-effect relationships
    \item Generate emergent strategic patterns
    \item Remain understandable for players and developers
    \item Produce varied and interesting game states
\end{itemize}

\subsection{Water Flow Simulation (Simple Model)}

Water is the fundamental resource connecting all game systems.

\subsubsection{Water Spreading Mechanism}

Each turn, water propagates through the grid using these rules:

\begin{enumerate}
    \item \textbf{Source Identification}
    \begin{itemize}
        \item River cells act as primary water sources
        \item Reservoir cells act as secondary water sources
        \item Each source has a water value (e.g., rivers = 10, reservoirs = 5--8)
    \end{itemize}
    
    \item \textbf{Adjacency Spreading}
    \begin{itemize}
        \item Water spreads to adjacent cells (4-directional or 8-directional)
        \item Water value decreases with distance from source
        \item Formula: $ \text{cell\_water} = \max(\text{neighbor\_water}) - 1 $
        \item Minimum water value is 0 (no water)
    \end{itemize}
    
    \item \textbf{Reservoir Amplification}
    \begin{itemize}
        \item Reservoir cells capture and store water
        \item Provide consistent water value to surrounding cells
        \item Boost water level in radius of 2--3 cells
        \item Prevent water value degradation in their vicinity
    \end{itemize}
    
    \item \textbf{Consumption Effects}
    \begin{itemize}
        \item Farm cells consume water each turn (reduces nearby water values)
        \item Multiple farms can deplete water sources
        \item Water consumption rate: 1--2 water units per farm per turn
    \end{itemize}
\end{enumerate}

\subsubsection{Water Meter Calculation}

The global water meter is calculated as:
\begin{equation}
\text{Water Meter} = \frac{\sum_{i=1}^{n} \text{cell\_water}_i}{n} \times \text{scaling\_factor}
\end{equation}

Where $n$ is the total number of cells, and scaling\_factor normalizes to 0--100 range.

\subsection{Crop and Farming System}

The agricultural system balances food production with water consumption.

\subsubsection{Crop Growth Cycle}

\begin{enumerate}
    \item \textbf{Planting Phase}
    \begin{itemize}
        \item Player places crop on land cell
        \item Crop starts at growth stage 0
        \item Requires adjacent water source (water value $\geq$ 3)
    \end{itemize}
    
    \item \textbf{Growth Progression}
    \begin{itemize}
        \item Each turn with sufficient water: growth stage $+1$
        \item Growth stages: 0 (planted) $\rightarrow$ 1 (sprouting) $\rightarrow$ 2 (growing) $\rightarrow$ 3 (mature)
        \item Insufficient water: growth stalls or regresses by 1 stage
        \item Total turns to maturity: 3--4 turns with consistent water
    \end{itemize}
    
    \item \textbf{Maturity and Harvest}
    \begin{itemize}
        \item Crops at stage 3 are harvestable
        \item Harvest action yields food points (5--10 points per harvest)
        \item After harvest, crop returns to stage 1 (faster regrowth)
    \end{itemize}
    
    \item \textbf{Water Dependency}
    \begin{itemize}
        \item Crops must have water value $\geq$ 3 in adjacent cells
        \item Distance from water affects growth rate:
        \begin{itemize}
            \item Adjacent to river/reservoir: 100\% growth rate
            \item 1 cell away: 75\% growth rate
            \item 2 cells away: 50\% growth rate
            \item 3+ cells away: 0\% growth rate (crop dies)
        \end{itemize}
    \end{itemize}
\end{enumerate}

\subsubsection{Farm Sustainability}

\begin{itemize}
    \item \textbf{Water Consumption Rate}: Each farm consumes 1--2 water units per turn
    \item \textbf{Overfarming Penalty}: Too many farms in one area deplete water sources
    \item \textbf{Optimal Density}: Approximately 1 farm per 3--4 water source cells
    \item \textbf{Food Meter Contribution}: Each mature farm contributes to global food meter
\end{itemize}

\subsection{Forest Ecosystem}

Forests provide long-term environmental benefits but require space and time investment.

\subsubsection{Forest Growth and Maturity}

\begin{enumerate}
    \item \textbf{Initial Planting}
    \begin{itemize}
        \item Young forest starts at maturity level 1
        \item Provides minimal oxygen production (1--2 units)
        \item Slight temperature reduction ($-0.5$ degrees)
    \end{itemize}
    
    \item \textbf{Maturation Process}
    \begin{itemize}
        \item Each turn, forest maturity increases: 1 $\rightarrow$ 2 $\rightarrow$ 3 (max)
        \item Maturation time: 3--5 turns to reach maximum effectiveness
        \item Maturity affects oxygen production and cooling power
    \end{itemize}
    
    \item \textbf{Mature Forest Benefits}
    \begin{itemize}
        \item Oxygen production: 5--8 units per turn
        \item Temperature reduction: $-1$ to $-2$ degrees per turn
        \item Ecosystem stability bonus
        \item May slightly increase local water retention
    \end{itemize}
\end{enumerate}

\subsubsection{Forest Density Effects}

Multiple adjacent forests create synergistic benefits:
\begin{itemize}
    \item \textbf{Small Forest} (1--2 trees): Standard benefits
    \item \textbf{Medium Forest} (3--5 adjacent trees): +20\% efficiency bonus
    \item \textbf{Large Forest} (6+ adjacent trees): +40\% efficiency bonus, creates "forest biome"
\end{itemize}

\subsubsection{Oxygen Production Calculation}

Global oxygen meter calculation:
\begin{equation}
\text{Oxygen Meter} = \sum_{i=1}^{m} (\text{forest\_maturity}_i \times \text{base\_oxygen})
\end{equation}

Where $m$ is the number of forest cells, and base\_oxygen = 2--3 units.

\subsection{Temperature System}

Temperature regulation is a central balancing challenge requiring careful management.

\subsubsection{Temperature Influencing Factors}

Temperature changes each turn based on multiple factors:

\begin{equation}
\Delta T = T_{\text{farming}} + T_{\text{industry}} - T_{\text{forest}} - T_{\text{water}}
\end{equation}

Where:
\begin{itemize}
    \item $ T_{\text{farming}} = \text{num\_farms} \times 0.3 $ (farming increases temp)
    \item $ T_{\text{industry}} = \text{num\_solar\_plants} \times 0.5 $ (industry increases temp)
    \item $ T_{\text{forest}} = \text{forest\_cooling\_power} $ (forests decrease temp)
    \item $ T_{\text{water}} = \text{water\_coverage} \times 0.2 $ (water moderates temp)
\end{itemize}

\subsubsection{Temperature Ranges and Effects}

\begin{itemize}
    \item \textbf{Optimal Range (40--60)}:
    \begin{itemize}
        \item Maximum crop growth efficiency
        \item No penalties to planetary systems
        \item Best oxygen production
    \end{itemize}
    
    \item \textbf{Moderate Range (25--40 or 60--75)}:
    \begin{itemize}
        \item Slightly reduced crop growth (75\% efficiency)
        \item Minor ecosystem stress
        \item Small stability penalty
    \end{itemize}
    
    \item \textbf{Extreme Cold (0--25)}:
    \begin{itemize}
        \item Severely reduced crop growth (25\% efficiency)
        \item Increased water freezing (reduced water availability)
        \item Major stability penalties
        \item Potential ecosystem collapse
    \end{itemize}
    
    \item \textbf{Extreme Heat (75--100)}:
    \begin{itemize}
        \item Crops wither faster (25\% efficiency)
        \item Increased water evaporation (reduced water meter)
        \item Forest fire risk (potential forest loss)
        \item Major stability penalties
        \item Potential ecosystem collapse
    \end{itemize}
\end{itemize}

\subsubsection{Temperature Victory Condition}

The goal is maintaining temperature in the safe range for sustained periods:
\begin{itemize}
    \item \textbf{Winning Strategy}: Keep temperature between 40--60 for 70\%+ of game
    \item \textbf{Penalty System}: Every turn outside optimal range accumulates damage
    \item \textbf{Recovery}: Returning to optimal range slowly heals accumulated damage
\end{itemize}

\subsection{Planet Stability Score}

The stability score synthesizes all metrics into a single health indicator.

\subsubsection{Stability Calculation Formula}

\begin{equation}
\text{Stability} = w_1 \cdot f(\text{water}) + w_2 \cdot f(\text{oxygen}) + w_3 \cdot f(\text{food}) + w_4 \cdot f(\text{temp})
\end{equation}

Where:
\begin{itemize}
    \item $ w_1, w_2, w_3, w_4 $ are weights (typically 0.25 each for balanced importance)
    \item $ f(\text{metric}) $ is a normalization function that scores each metric:
\end{itemize}

\begin{equation}
f(\text{metric}) = 
\begin{cases}
1.0 & \text{if metric in optimal range (50--80)} \\
0.5 + 0.5 \cdot \frac{\text{metric}}{50} & \text{if metric below 50} \\
1.0 - 0.5 \cdot \frac{\text{metric} - 80}{20} & \text{if metric above 80}
\end{cases}
\end{equation}

\begin{itemize}
    \item For temperature, $ f(\text{temp}) $ uses the optimal range of 40--60
\end{itemize}

\subsubsection{Stability Thresholds}

\begin{itemize}
    \item \textbf{High Stability (0.8--1.0)}:
    \begin{itemize}
        \item All resources balanced
        \item Temperature in optimal range
        \item Sustainable ecosystem
        \item Bonus points awarded
    \end{itemize}
    
    \item \textbf{Moderate Stability (0.5--0.8)}:
    \begin{itemize}
        \item Some resources stressed
        \item Temperature slightly off-target
        \item Ecosystem functional but vulnerable
        \item No bonuses or penalties
    \end{itemize}
    
    \item \textbf{Low Stability (0.2--0.5)}:
    \begin{itemize}
        \item One or more resources critical
        \item Temperature approaching extremes
        \item Ecosystem under severe stress
        \item Penalty points accumulate
    \end{itemize}
    
    \item \textbf{Critical Instability (0.0--0.2)}:
    \begin{itemize}
        \item Resource collapse imminent or occurring
        \item Extreme temperature conditions
        \item Ecosystem collapse
        \item Rapid accumulation of penalties
        \item Potential game over condition
    \end{itemize}
\end{itemize}

\subsubsection{Win/Loss Conditions Based on Stability}

\begin{itemize}
    \item \textbf{Victory}: Maintain average stability $> 0.7$ for 50+ turns
    \item \textbf{Defeat}: Stability drops below 0.15 for 3 consecutive turns
    \item \textbf{Score-Based Win}: Highest stability score after fixed number of turns (100--150)
\end{itemize}

% ================================================================
% SECTION 4: AI AGENT ARCHITECTURE
% ================================================================
\section{How to Build the AI Agents}

This section provides a comprehensive guide to implementing intelligent agents capable of playing the game strategically.

\subsection{Step 1: State Representation}

The AI agent must have complete visibility into the game state to make informed decisions.

\subsubsection{State Components}

The complete game state includes the following information:

\begin{enumerate}
    \item \textbf{Grid Layout}
    \begin{itemize}
        \item 2D array representing all cells (25 $\times$ 18)
        \item Each cell contains:
        \begin{itemize}
            \item Terrain type (land, river, farm, forest, reservoir, solar plant)
            \item Owner (Agent A, Agent B, or neutral)
            \item Additional properties:
            \begin{itemize}
                \item Water value (0--10)
                \item Crop growth stage (0--3) if farm
                \item Forest maturity (1--3) if forest
                \item Gate state (open/closed) if applicable
            \end{itemize}
        \end{itemize}
    \end{itemize}
    
    \item \textbf{Global Planetary Meters}
    \begin{itemize}
        \item Water level (0--100)
        \item Oxygen level (0--100)
        \item Food supply (0--100)
        \item Temperature (0--100)
        \item Planet stability score (0.0--1.0)
    \end{itemize}
    
    \item \textbf{Game Progress Information}
    \begin{itemize}
        \item Current turn number
        \item Total turns elapsed
        \item Turns remaining (if using turn limit)
    \end{itemize}
    
    \item \textbf{Player Information}
    \begin{itemize}
        \item Own score/points
        \item Opponent score/points (optional, depending on game mode)
        \item Available eco points for actions
        \item Number of tiles controlled by each player
    \end{itemize}
    
    \item \textbf{Historical Data (Optional but Useful)}
    \begin{itemize}
        \item Last 5--10 actions taken by each player
        \item Meter trend data (improving/declining)
        \item Stability history
    \end{itemize}
\end{enumerate}

\subsubsection{State Encoding for AI}

The state can be encoded in multiple formats depending on the AI algorithm:

\begin{itemize}
    \item \textbf{Array/Matrix Representation}:
    \begin{itemize}
        \item Grid as multi-channel 2D array (one channel per property)
        \item Meters as separate vector [water, oxygen, food, temp, stability]
        \item Efficient for neural network approaches
    \end{itemize}
    
    \item \textbf{Object-Oriented Representation}:
    \begin{itemize}
        \item State as a class with methods to query specific information
        \item Easy to work with in traditional programming
        \item Better for tree search algorithms
    \end{itemize}
    
    \item \textbf{Feature Vector Representation}:
    \begin{itemize}
        \item Extract key features: \# forests, \# farms, avg water, etc.
        \item Compressed representation for faster computation
        \item Useful for machine learning classifiers
    \end{itemize}
\end{itemize}

\subsection{Step 2: Move Generator (Action Space)}

The move generator identifies all legal actions available in the current state.

\subsubsection{Valid Action Detection}

For each action type, determine valid locations:

\begin{enumerate}
    \item \textbf{Build Canal}
    \begin{itemize}
        \item Valid cells: Land cells not owned by opponent
        \item Constraint: Must have eco points $\geq$ 3
        \item Constraint: Cannot be adjacent to existing canal (optional rule)
    \end{itemize}
    
    \item \textbf{Place Reservoir}
    \begin{itemize}
        \item Valid cells: Land or river cells
        \item Constraint: Must be near existing water source (within 2--3 cells)
        \item Constraint: Eco points $\geq$ 6
    \end{itemize}
    
    \item \textbf{Open/Close Gate}
    \begin{itemize}
        \item Valid cells: Existing gate tiles only
        \item No location choice, just toggle action
        \item Low eco point cost (1--2)
    \end{itemize}
    
    \item \textbf{Plant Forest}
    \begin{itemize}
        \item Valid cells: Empty land cells
        \item Constraint: Cannot be adjacent to farm (optional spacing rule)
        \item Constraint: Eco points $\geq$ 4
    \end{itemize}
    
    \item \textbf{Clear Forest}
    \begin{itemize}
        \item Valid cells: Any forest cell
        \item Preference: Clear opponent's forests or old/inefficient forests
        \item Environmental penalty applied
    \end{itemize}
    
    \item \textbf{Plant Crop}
    \begin{itemize}
        \item Valid cells: Land cells with water value $\geq$ 3 nearby
        \item Constraint: Eco points $\geq$ 3
        \item Constraint: Not already a farm
    \end{itemize}
    
    \item \textbf{Harvest Crop}
    \begin{itemize}
        \item Valid cells: Farm cells with growth stage = 3
        \item Constraint: Must have sufficient water to complete harvest
    \end{itemize}
    
    \item \textbf{Build Solar Plant}
    \begin{itemize}
        \item Valid cells: Empty land cells
        \item Constraint: Eco points $\geq$ 8
        \item Consideration: Strategic placement for energy distribution
    \end{itemize}
\end{enumerate}

\subsubsection{Action Pruning for Efficiency}

To reduce computational load, implement action pruning:

\begin{itemize}
    \item \textbf{Relevance Filtering}: Remove actions that don't address current needs
    \item \textbf{Proximity Grouping}: Only consider actions in "active" regions of map
    \item \textbf{Resource Gating}: Skip expensive actions when eco points are low
    \item \textbf{Heuristic Pre-screening}: Use fast heuristics to eliminate obviously bad moves
\end{itemize}

\subsection{Step 3: Evaluation Function (The Key to Intelligence)}

The evaluation function is the heart of the AI, scoring how good a particular game state is.

\subsubsection{Core Evaluation Formula}

A comprehensive evaluation function combines multiple factors:

\begin{equation}
\text{score}(\text{state}) = \sum_{i} w_i \cdot f_i(\text{state})
\end{equation}

Where each $f_i$ is a component evaluating a specific aspect:

\begin{enumerate}
    \item \textbf{Point Differential (Competitive Component)}
    \begin{equation}
    f_1 = (\text{my\_points} - \text{opponent\_points})
    \end{equation}
    Weight: $w_1 = 5.0$ (highly important in competitive mode)
    
    \item \textbf{Planet Stability (Cooperative/Survival Component)}
    \begin{equation}
    f_2 = \text{planet\_stability}
    \end{equation}
    Weight: $w_2 = 10.0$ (critical for long-term survival)
    
    \item \textbf{Resource Levels (Individual Meters)}
    \begin{equation}
    f_3 = \text{water\_score} + \text{oxygen\_score} + \text{food\_score}
    \end{equation}
    
    Where each resource score uses a target-based function:
    \begin{equation}
    \text{resource\_score} = 
    \begin{cases}
    \frac{\text{level}}{50} & \text{if level} < 50 \text{ (penalty for low)} \\
    1.0 & \text{if } 50 \leq \text{level} \leq 80 \text{ (optimal)} \\
    1.0 - 0.5 \cdot \frac{\text{level} - 80}{20} & \text{if level} > 80 \text{ (diminishing returns)}
    \end{cases}
    \end{equation}
    Weight: $w_3 = 8.0$
    
    \item \textbf{Temperature Optimality}
    \begin{equation}
    f_4 = 1.0 - \frac{|\text{temp} - 50|}{50}
    \end{equation}
    This gives a score of 1.0 when temp = 50, and decreases linearly toward extremes.
    
    Weight: $w_4 = 12.0$ (temperature is critical)
    
    \item \textbf{Resource Collapse Penalty}
    \begin{equation}
    f_5 = -\sum_{r \in \text{resources}} 
    \begin{cases}
    100 & \text{if } r < 20 \text{ (critical)} \\
    20 & \text{if } 20 \leq r < 40 \text{ (danger)} \\
    0 & \text{otherwise}
    \end{cases}
    \end{equation}
    Weight: $w_5 = 1.0$ (already large penalties)
    
    \item \textbf{Strategic Asset Value}
    \begin{equation}
    f_6 = \text{num\_forests} \times 2 + \text{num\_farms} \times 3 + \text{num\_reservoirs} \times 4
    \end{equation}
    Weight: $w_6 = 1.5$ (encourages infrastructure building)
    
    \item \textbf{Territory Control (Optional, for Competitive Mode)}
    \begin{equation}
    f_7 = \text{my\_controlled\_cells} - \text{opponent\_controlled\_cells}
    \end{equation}
    Weight: $w_7 = 2.0$
\end{enumerate}

\subsubsection{Evaluation Function Tuning}

Weights can be adjusted based on game phase:

\begin{itemize}
    \item \textbf{Early Game} (turns 1--30):
    \begin{itemize}
        \item Increase weight on infrastructure ($w_6$)
        \item Decrease weight on immediate resources
        \item Focus on long-term positioning
    \end{itemize}
    
    \item \textbf{Mid Game} (turns 31--70):
    \begin{itemize}
        \item Balanced weights
        \item Emphasis on stability and resource balance
    \end{itemize}
    
    \item \textbf{Late Game} (turns 71+):
    \begin{itemize}
        \item Increase weight on point differential ($w_1$)
        \item Increase temperature management ($w_4$)
        \item Focus on maintaining leads or catching up
    \end{itemize}
\end{itemize}

\subsection{Three Levels of AI Implementation}

Implementing multiple AI difficulty levels allows for algorithm comparison and player choice.

\subsubsection{AI Level 1: Greedy Heuristic (Fast, Good Baseline)}

\textbf{Algorithm Description:}

The greedy AI selects the action that immediately maximizes the evaluation function.

\textbf{Pseudocode:}
\begin{lstlisting}
function greedy_ai(state):
    best_action = null
    best_score = -infinity
    
    for each valid_action in get_valid_actions(state):
        # Simulate one step ahead
        next_state = simulate_action(state, valid_action)
        score = evaluate(next_state)
        
        if score > best_score:
            best_score = score
            best_action = valid_action
    
    return best_action
\end{lstlisting}

\textbf{Characteristics:}
\begin{itemize}
    \item \textbf{Computational Complexity}: $O(n)$ where $n$ is number of valid actions
    \item \textbf{Strength}: Fast, makes locally optimal decisions
    \item \textbf{Weakness}: No lookahead, can be tricked by short-term gains
    \item \textbf{Best For}: Beginner difficulty, baseline comparison
\end{itemize}

\textbf{Expected Performance}:
\begin{itemize}
    \item Responds well to immediate threats (e.g., low food)
    \item Builds infrastructure when resources are stable
    \item Vulnerable to strategic traps
    \item Average game score: 60--70\% of optimal
\end{itemize}

\subsubsection{AI Level 2: Minimax with Alpha-Beta Pruning (Impressive for Lab)}

\textbf{Algorithm Description:}

Minimax assumes optimal opponent play and searches the game tree to depth 2--3, alternating between maximizing own score and minimizing opponent's advantage.

\textbf{Pseudocode:}
\begin{lstlisting}
function minimax_ai(state, depth):
    best_action = null
    best_value = -infinity
    
    for each valid_action in get_valid_actions(state):
        next_state = simulate_action(state, valid_action)
        value = minimax_value(next_state, depth - 1, -infinity, +infinity, false)
        
        if value > best_value:
            best_value = value
            best_action = valid_action
    
    return best_action

function minimax_value(state, depth, alpha, beta, is_maximizing):
    if depth == 0 or is_terminal(state):
        return evaluate(state)
    
    if is_maximizing:
        max_value = -infinity
        for each action in get_valid_actions(state):
            next_state = simulate_action(state, action)
            value = minimax_value(next_state, depth - 1, alpha, beta, false)
            max_value = max(max_value, value)
            alpha = max(alpha, value)
            if beta <= alpha:
                break  # Beta cutoff
        return max_value
    else:
        min_value = +infinity
        for each action in get_opponent_actions(state):
            next_state = simulate_action(state, action)
            value = minimax_value(next_state, depth - 1, alpha, beta, true)
            min_value = min(min_value, value)
            beta = min(beta, value)
            if beta <= alpha:
                break  # Alpha cutoff
        return min_value
\end{lstlisting}

\textbf{Characteristics:}
\begin{itemize}
    \item \textbf{Computational Complexity}: $O(b^d)$ where $b$ is branching factor, $d$ is depth
    \item \textbf{With Alpha-Beta}: Can reduce to $O(b^{d/2})$ with good move ordering
    \item \textbf{Strength}: Anticipates opponent moves, strategic planning
    \item \textbf{Weakness}: Computationally expensive at depth > 3
    \item \textbf{Best For}: Intermediate to advanced difficulty, demonstrates game theory
\end{itemize}

\textbf{Optimization Techniques:}
\begin{itemize}
    \item \textbf{Move Ordering}: Evaluate promising moves first to maximize pruning
    \item \textbf{Transposition Tables}: Cache evaluated states to avoid recomputation
    \item \textbf{Iterative Deepening}: Start with depth 1, increase gradually with time
    \item \textbf{Quiescence Search}: Extend search at critical decision points
\end{itemize}

\textbf{Expected Performance}:
\begin{itemize}
    \item Anticipates 2--3 turns ahead
    \item Counters opponent strategies effectively
    \item Makes sacrificial plays for long-term advantage
    \item Average game score: 80--90\% of optimal
\end{itemize}

\subsubsection{AI Level 3: Monte Carlo Tree Search (Advanced, Impressive)}

\textbf{Algorithm Description:}

MCTS samples many random game continuations (rollouts) from each possible action and selects the action with the best average outcome.

\textbf{Pseudocode:}
\begin{lstlisting}
function mcts_ai(state, num_rollouts):
    action_scores = {}
    
    for each valid_action in get_valid_actions(state):
        total_score = 0
        
        # Perform multiple rollouts
        for i = 1 to num_rollouts:
            simulated_state = simulate_action(state, valid_action)
            
            # Random or semi-intelligent rollout
            rollout_score = perform_rollout(simulated_state, rollout_depth)
            total_score += rollout_score
        
        # Average score for this action
        action_scores[valid_action] = total_score / num_rollouts
    
    # Select action with best average outcome
    return argmax(action_scores)

function perform_rollout(state, depth):
    current_state = state
    
    for turn = 1 to depth:
        if is_terminal(current_state):
            break
        
        # Select random or heuristic-guided action
        action = select_rollout_action(current_state)
        current_state = simulate_action(current_state, action)
    
    return evaluate(current_state)

function select_rollout_action(state):
    # Option 1: Pure random
    return random_choice(get_valid_actions(state))
    
    # Option 2: Semi-intelligent (better results)
    actions = get_valid_actions(state)
    action_weights = [quick_heuristic(state, a) for a in actions]
    return weighted_random_choice(actions, action_weights)
\end{lstlisting}

\textbf{Characteristics:}
\begin{itemize}
    \item \textbf{Computational Complexity}: $O(n \times r \times d)$ where $n$ = actions, $r$ = rollouts, $d$ = depth
    \item \textbf{Strength}: Handles uncertainty well, discovers non-obvious strategies
    \item \textbf{Weakness}: Requires many rollouts for accuracy, can be computationally intensive
    \item \textbf{Best For}: Advanced difficulty, demonstrates modern AI techniques
\end{itemize}

\textbf{Enhancement Options:}
\begin{itemize}
    \item \textbf{UCB1 Selection}: Use Upper Confidence Bound to balance exploration/exploitation
    \item \textbf{Progressive Widening}: Gradually increase rollout depth as confidence grows
    \item \textbf{Domain Knowledge}: Guide rollouts with heuristics instead of pure random
    \item \textbf{Parallelization}: Run rollouts in parallel for speed
\end{itemize}

\textbf{Expected Performance}:
\begin{itemize}
    \item Adapts to unexpected opponent strategies
    \item Finds creative solutions to resource problems
    \item Excellent at balancing multiple objectives
    \item Average game score: 90--95\% of optimal (with sufficient rollouts)
\end{itemize}

\subsection{AI Comparison Metrics for Laboratory Analysis}

To demonstrate AI effectiveness in an academic setting, measure:

\begin{enumerate}
    \item \textbf{Win Rate}: Percentage of games won against each other AI
    \item \textbf{Average Stability Score}: Long-term planetary health maintenance
    \item \textbf{Resource Efficiency}: Points earned per resource consumed
    \item \textbf{Decision Time}: Computational time per move
    \item \textbf{Adaptability}: Performance against different opponent strategies
    \item \textbf{Learning Curve}: Improvement over successive games (if using learning)
\end{enumerate}

% ================================================================
% SECTION 5: VISUAL PRESENTATION
% ================================================================
\section{Making It Visually Beautiful (Without Unity)}

Creating a professional-looking game does not require expensive engines like Unity. Python with Pygame and quality assets can produce impressive results.

\subsection{Recommended Technology Stack}

\subsubsection{Core Framework: Python + Pygame}

\textbf{Why Python + Pygame?}
\begin{itemize}
    \item \textbf{Simplicity}: Easy to learn and implement
    \item \textbf{Rapid Development}: Fast iteration on game logic and visuals
    \item \textbf{Cross-Platform}: Works on Windows, Mac, Linux
    \item \textbf{AI Integration}: Python has excellent AI libraries
    \item \textbf{Free and Open Source}: No licensing costs
\end{itemize}

\textbf{Installation}:
\begin{lstlisting}[language=bash]
pip install pygame
pip install numpy  # For efficient grid operations
\end{lstlisting}

\subsubsection{Asset Sources (Free, High-Quality)}

To avoid the "colored squares" aesthetic, use professional tile sets:

\begin{enumerate}
    \item \textbf{Kenney.nl} (Highly Recommended)
    \begin{itemize}
        \item Website: \url{https://kenney.nl/assets}
        \item Content: Thousands of free 2D game assets
        \item Styles: Pixel art, flat design, isometric
        \item License: CC0 (public domain)
        \item Recommended packs:
        \begin{itemize}
            \item "Topdown Tanks" (can repurpose tiles)
            \item "Isometric Prototype Tiles"
            \item "RPG Urban Pack"
            \item "Nature Kit"
        \end{itemize}
    \end{itemize}
    
    \item \textbf{OpenGameArt.org}
    \begin{itemize}
        \item Website: \url{https://opengameart.org/}
        \item Content: Community-contributed game art
        \item Search terms: "terrain tiles", "farm sprites", "water animation"
        \item License: Varies (check individual assets)
        \item Good for: Specific tiles, character sprites, UI elements
    \end{itemize}
    
    \item \textbf{Itch.io Asset Packs}
    \begin{itemize}
        \item Website: \url{https://itch.io/game-assets/free}
        \item Content: Many free and paid high-quality assets
        \item Filter: Use "free" tag
        \item Good for: Cohesive themed packs
    \end{itemize}
\end{enumerate}

\subsection{Visual Style Checklist}

To create a polished, professional appearance:

\subsubsection{Tilemap Implementation}

\begin{itemize}
    \item \textbf{Tile Sizes}: Use consistent tile dimensions (32×32, 64×64, or 128×128 pixels)
    \item \textbf{Terrain Variety}:
    \begin{itemize}
        \item Grass: Multiple variants (light, dark, with flowers)
        \item Water: Animated tiles with ripples
        \item Trees: Different species and sizes
        \item Farms: Show different crop types and growth stages
        \item Buildings: Distinct sprites for reservoir, solar plant, etc.
    \end{itemize}
    \item \textbf{Tile Transitions}: Use edge tiles for smooth transitions between terrain types
    \item \textbf{Auto-Tiling}: Implement auto-tiling algorithm for natural-looking terrain borders
\end{itemize}

\subsubsection{Water Animation}

Even simple 2-frame animation creates life:

\begin{lstlisting}[language=Python]
# Pseudo-code for water animation
water_frames = [load_image('water_1.png'), load_image('water_2.png')]
current_frame = 0
animation_timer = 0
ANIMATION_SPEED = 30  # Change frame every 30 ticks

def update():
    global animation_timer, current_frame
    animation_timer += 1
    if animation_timer >= ANIMATION_SPEED:
        animation_timer = 0
        current_frame = (current_frame + 1) % len(water_frames)

def draw_water(x, y):
    screen.blit(water_frames[current_frame], (x, y))
\end{lstlisting}

\subsubsection{Resource Visualization Effects}

Make resources feel dynamic:

\begin{itemize}
    \item \textbf{Glow Effects}: Add pulsing glow to crystals/resources
    \begin{lstlisting}[language=Python]
# Pseudo-code for pulsing glow
glow_intensity = 128 + 127 * sin(time * 2)  # Oscillates 0-255
draw_circle_with_alpha(position, radius, color, glow_intensity)
    \end{lstlisting}
    
    \item \textbf{Particle Systems}: Small particles for visual feedback
    \begin{itemize}
        \item Harvest action: Floating "+Food" text with upward motion
        \item Forest planting: Green sparkles
        \item Building construction: Dust clouds
    \end{itemize}
\end{itemize}

\subsubsection{User Interface Panel}

Design a clean, informative UI on the right side:

\begin{itemize}
    \item \textbf{Meter Bars}: Visual bars for Water, Oxygen, Food, Temperature
    \begin{lstlisting}[language=Python]
def draw_meter(x, y, label, value, max_value, color):
    # Background bar
    draw_rect(x, y, 200, 30, GRAY)
    # Filled bar
    fill_width = (value / max_value) * 200
    draw_rect(x, y, fill_width, 30, color)
    # Label and value text
    draw_text(label + ": " + str(value), x, y - 20, BLACK)
    \end{lstlisting}
    
    \item \textbf{Score Display}: Large, readable font showing points
    \item \textbf{Turn Counter}: Current turn / max turns
    \item \textbf{Action Log}: Text box showing last 3--5 actions
    \begin{itemize}
        \item "Agent A built canal at (12, 8)"
        \item "Agent B harvested crop (+8 food)"
    \end{itemize}
    \item \textbf{Planet Status}: Large icon showing stability (happy/neutral/sad planet)
\end{itemize}

\subsection{Professional Visual Effects (Easy to Implement)}

\subsubsection{Smooth Tile-to-Tile Movement}

Instead of instant placement, interpolate movement:

\begin{lstlisting}[language=Python]
class AnimatedSprite:
    def __init__(self, start_pos, end_pos, duration):
        self.start = start_pos
        self.end = end_pos
        self.duration = duration
        self.elapsed = 0
    
    def update(self, dt):
        self.elapsed += dt
        if self.elapsed >= self.duration:
            self.elapsed = self.duration
    
    def get_position(self):
        t = self.elapsed / self.duration  # 0 to 1
        # Linear interpolation
        x = self.start[0] + (self.end[0] - self.start[0]) * t
        y = self.start[1] + (self.end[1] - self.start[1]) * t
        return (x, y)
\end{lstlisting}

\subsubsection{Particle Pop Effects}

When harvesting or building, show floating feedback:

\begin{lstlisting}[language=Python]
class FloatingText:
    def __init__(self, text, position, color):
        self.text = text
        self.x, self.y = position
        self.color = color
        self.lifetime = 60  # frames
        self.age = 0
    
    def update(self):
        self.age += 1
        self.y -= 1  # Float upward
        self.alpha = 255 * (1 - self.age / self.lifetime)
    
    def draw(self, screen):
        if self.age < self.lifetime:
            draw_text_with_alpha(screen, self.text, self.x, self.y, 
                               self.color, self.alpha)
    
    def is_alive(self):
        return self.age < self.lifetime

# Usage
floating_texts = []
# When harvesting
floating_texts.append(FloatingText("+8 Food", harvest_pos, GREEN))
\end{lstlisting}

\subsubsection{Day/Night Cycle (Ambient Atmosphere)}

Subtle lighting changes add depth:

\begin{lstlisting}[language=Python]
def get_time_of_day_tint(turn_number):
    # Cycle every 20 turns: day -> dusk -> night -> dawn
    phase = (turn_number % 20) / 20.0
    
    if phase < 0.25:  # Day
        return (255, 255, 255)
    elif phase < 0.5:  # Dusk
        return (255, 200, 150)
    elif phase < 0.75:  # Night
        return (100, 100, 150)
    else:  # Dawn
        return (200, 150, 200)

def apply_tint(screen, tint):
    overlay = pygame.Surface(screen.get_size())
    overlay.fill(tint)
    screen.blit(overlay, (0, 0), special_flags=pygame.BLEND_MULT)
\end{lstlisting}

\subsubsection{Camera Shake (For Dramatic Events)}

When critical events occur (resource collapse, major construction):

\begin{lstlisting}[language=Python]
class CameraShake:
    def __init__(self):
        self.shake_amount = 0
        self.shake_duration = 0
    
    def trigger(self, amount, duration):
        self.shake_amount = amount
        self.shake_duration = duration
    
    def update(self):
        if self.shake_duration > 0:
            self.shake_duration -= 1
            return (random.randint(-self.shake_amount, self.shake_amount),
                   random.randint(-self.shake_amount, self.shake_amount))
        return (0, 0)

# Usage
camera_shake = CameraShake()
# When temperature reaches extreme
if temperature > 90:
    camera_shake.trigger(amount=5, duration=20)

# In render loop
offset = camera_shake.update()
screen.blit(game_surface, offset)
\end{lstlisting}

\subsection{Performance Optimization Tips}

\begin{itemize}
    \item \textbf{Dirty Rectangle Rendering}: Only redraw changed portions of screen
    \item \textbf{Sprite Sheets}: Load all tiles from single image to reduce I/O
    \item \textbf{Surface Caching}: Pre-render complex UI elements
    \item \textbf{Event-Driven Updates}: Only recalculate when state changes
    \item \textbf{Target 60 FPS}: Cap frame rate for consistent performance
\end{itemize}

% ================================================================
% SECTION 6: PROJECT STRUCTURE
% ================================================================
\section{Recommended Project Structure}

A well-organized project structure is essential for maintainability, debugging, and presenting professional work to instructors.

\subsection{Directory Layout}

\begin{lstlisting}
world_balance_architects/
|
+-- main.py                    # Entry point, main game loop
+-- config.py                  # Global configuration and constants
+-- requirements.txt           # Python dependencies
+-- README.md                  # Project documentation
|
+-- assets/                    # All visual and audio resources
|   |
|   +-- tiles/                 # Tile images
|   |   +-- grass.png
|   |   +-- water_1.png
|   |   +-- water_2.png
|   |   +-- forest.png
|   |   +-- farm.png
|   |   +-- reservoir.png
|   |   +-- solar_plant.png
|   |
|   +-- ui/                    # User interface elements
|       +-- panel_background.png
|       +-- button.png
|       +-- meter_bar.png
|       +-- icons/
|           +-- water_icon.png
|           +-- food_icon.png
|           +-- oxygen_icon.png
|           +-- temp_icon.png
|
+-- engine/                    # Core game logic
|   |
|   +-- world.py               # Grid representation and management
|   |                          # - World class
|   |                          # - Cell class
|   |                          # - Grid initialization
|   |
|   +-- simulate.py            # Simulation rules and updates
|   |                          # - Water flow simulation
|   |                          # - Crop growth logic
|   |                          # - Forest maturation
|   |                          # - Temperature calculation
|   |                          # - Stability score computation
|   |
|   +-- actions.py             # Action definitions and validation
|                              # - Action base class
|                              # - BuildCanal, PlaceReservoir, etc.
|                              # - validate_action()
|                              # - apply_action()
|
+-- ai/                        # AI agent implementations
|   |
|   +-- base.py                # Abstract base AI class
|   |                          # - AIAgent interface
|   |                          # - Common helper methods
|   |
|   +-- greedy.py              # Greedy heuristic AI
|   |                          # - GreedyAI class
|   |
|   +-- minimax.py             # Minimax algorithm
|   |                          # - MinimaxAI class
|   |                          # - Alpha-beta pruning
|   |
|   +-- mcts.py                # Monte Carlo Tree Search
|   |                          # - MCTSAI class
|   |                          # - Rollout implementation
|   |
|   +-- eval.py                # Evaluation functions
|                              # - evaluate_state()
|                              # - Component scoring functions
|
+-- render/                    # Visualization and graphics
    |
    +-- renderer.py            # Main rendering system
    |                          # - Renderer class
    |                          # - draw_grid()
    |                          # - draw_ui()
    |                          # - Asset loading
    |
    +-- animations.py          # Animation systems
                               # - AnimatedSprite class
                               # - FloatingText class
                               # - ParticleSystem class
                               # - CameraShake class
\end{lstlisting}

\subsection{File Descriptions and Responsibilities}

\subsubsection{main.py}

Entry point for the application. Handles initialization and main game loop.

\begin{lstlisting}[language=Python]
# main.py structure
import pygame
from engine.world import World
from engine.simulate import simulate_turn
from ai.greedy import GreedyAI
from ai.minimax import MinimaxAI
from render.renderer import Renderer

def main():
    # Initialize Pygame
    pygame.init()
    screen = pygame.display.set_mode((1280, 720))
    clock = pygame.time.Clock()
    
    # Create game world
    world = World(width=25, height=18)
    
    # Create AI agents
    agent_a = MinimaxAI(depth=2)
    agent_b = GreedyAI()
    
    # Create renderer
    renderer = Renderer(screen)
    
    # Main game loop
    running = True
    while running:
        # Event handling
        for event in pygame.event.get():
            if event.type == pygame.QUIT:
                running = False
        
        # Agent A turn
        action_a = agent_a.choose_action(world)
        world.apply_action(action_a)
        simulate_turn(world)
        
        # Agent B turn
        action_b = agent_b.choose_action(world)
        world.apply_action(action_b)
        simulate_turn(world)
        
        # Render
        renderer.draw(world)
        pygame.display.flip()
        clock.tick(60)
    
    pygame.quit()

if __name__ == "__main__":
    main()
\end{lstlisting}

\subsubsection{config.py}

Centralized configuration for easy tuning.

\begin{lstlisting}[language=Python]
# config.py structure
# Grid settings
GRID_WIDTH = 25
GRID_HEIGHT = 18
TILE_SIZE = 48

# Game settings
MAX_TURNS = 150
TARGET_TEMPERATURE = 50
TEMPERATURE_TOLERANCE = 10

# Action costs
ACTION_COSTS = {
    'build_canal': 3,
    'place_reservoir': 6,
    'toggle_gate': 1,
    'plant_forest': 4,
    'clear_forest': 2,
    'plant_crop': 3,
    'harvest_crop': 1,
    'build_solar': 8
}

# Simulation parameters
WATER_FLOW_RATE = 1.5
CROP_GROWTH_RATE = 1.0
FOREST_MATURATION_TURNS = 4
TEMPERATURE_CHANGE_RATE = 0.5

# Evaluation function weights
EVAL_WEIGHTS = {
    'point_differential': 5.0,
    'stability': 10.0,
    'resources': 8.0,
    'temperature': 12.0,
    'collapse_penalty': 1.0,
    'assets': 1.5,
    'territory': 2.0
}

# Visual settings
SCREEN_WIDTH = 1280
SCREEN_HEIGHT = 720
FPS = 60
ANIMATION_SPEED = 30
\end{lstlisting}

\subsubsection{engine/world.py}

Manages the grid and world state.

\begin{lstlisting}[language=Python]
# engine/world.py structure
class Cell:
    def __init__(self, x, y, terrain_type='land'):
        self.x = x
        self.y = y
        self.terrain_type = terrain_type
        self.owner = None
        self.water_value = 0
        self.crop_growth_stage = 0
        self.forest_maturity = 0
        self.gate_open = False

class World:
    def __init__(self, width, height):
        self.width = width
        self.height = height
        self.grid = [[Cell(x, y) for y in range(height)] 
                     for x in range(width)]
        self.meters = {
            'water': 50,
            'oxygen': 50,
            'food': 50,
            'temperature': 50
        }
        self.turn = 0
        self.stability = 0.5
    
    def get_cell(self, x, y):
        if 0 <= x < self.width and 0 <= y < self.height:
            return self.grid[x][y]
        return None
    
    def get_neighbors(self, x, y, distance=1):
        neighbors = []
        for dx in range(-distance, distance + 1):
            for dy in range(-distance, distance + 1):
                if dx == 0 and dy == 0:
                    continue
                cell = self.get_cell(x + dx, y + dy)
                if cell:
                    neighbors.append(cell)
        return neighbors
    
    def apply_action(self, action):
        # Apply action to world
        action.execute(self)
\end{lstlisting}

\subsubsection{engine/simulate.py}

Implements all simulation rules.

\begin{lstlisting}[language=Python]
# engine/simulate.py structure
def simulate_turn(world):
    simulate_water_flow(world)
    simulate_crop_growth(world)
    simulate_forest_growth(world)
    calculate_temperature(world)
    update_meters(world)
    calculate_stability(world)
    world.turn += 1

def simulate_water_flow(world):
    # Water spreading algorithm
    for x in range(world.width):
        for y in range(world.height):
            cell = world.get_cell(x, y)
            if cell.terrain_type == 'river':
                cell.water_value = 10
            # Additional flow logic...

def simulate_crop_growth(world):
    # Crop maturation logic
    pass

def simulate_forest_growth(world):
    # Forest maturation logic
    pass

def calculate_temperature(world):
    # Temperature calculation based on factors
    pass

def update_meters(world):
    # Update global resource meters
    pass

def calculate_stability(world):
    # Compute planet stability score
    pass
\end{lstlisting}

\subsubsection{ai/base.py}

Abstract interface for all AI agents.

\begin{lstlisting}[language=Python]
# ai/base.py structure
from abc import ABC, abstractmethod

class AIAgent(ABC):
    def __init__(self, player_id):
        self.player_id = player_id
    
    @abstractmethod
    def choose_action(self, world):
        """
        Given the current world state, return the best action.
        """
        pass
    
    def get_valid_actions(self, world):
        """
        Returns list of all valid actions for current state.
        """
        actions = []
        # Enumerate all possible actions
        return actions
\end{lstlisting}

\subsubsection{render/renderer.py}

Handles all visual rendering.

\begin{lstlisting}[language=Python]
# render/renderer.py structure
import pygame

class Renderer:
    def __init__(self, screen):
        self.screen = screen
        self.assets = self.load_assets()
    
    def load_assets(self):
        """Load all images and sprites."""
        assets = {
            'grass': pygame.image.load('assets/tiles/grass.png'),
            'water': [
                pygame.image.load('assets/tiles/water_1.png'),
                pygame.image.load('assets/tiles/water_2.png')
            ],
            # ... load all other assets
        }
        return assets
    
    def draw(self, world):
        """Main rendering method."""
        self.screen.fill((200, 220, 255))  # Sky blue background
        self.draw_grid(world)
        self.draw_ui(world)
    
    def draw_grid(self, world):
        """Render the game grid."""
        for x in range(world.width):
            for y in range(world.height):
                cell = world.get_cell(x, y)
                self.draw_cell(cell)
    
    def draw_cell(self, cell):
        """Render individual cell."""
        sprite = self.assets.get(cell.terrain_type)
        pos = (cell.x * TILE_SIZE, cell.y * TILE_SIZE)
        self.screen.blit(sprite, pos)
    
    def draw_ui(self, world):
        """Render UI panel with meters and info."""
        self.draw_meter(950, 50, 'Water', world.meters['water'])
        self.draw_meter(950, 100, 'Oxygen', world.meters['oxygen'])
        # ... draw all meters
    
    def draw_meter(self, x, y, label, value):
        """Draw individual resource meter."""
        # Meter rendering code
        pass
\end{lstlisting}

\subsection{Development Workflow}

\subsubsection{Recommended Implementation Order}

\begin{enumerate}
    \item \textbf{Phase 1: Core Engine} (Week 1)
    \begin{itemize}
        \item Implement World and Cell classes
        \item Create basic grid representation
        \item Implement simple rendering (colored squares okay for now)
        \item Test grid creation and display
    \end{itemize}
    
    \item \textbf{Phase 2: Actions and Simulation} (Week 2)
    \begin{itemize}
        \item Implement all action types
        \item Add basic simulation rules (water flow, crop growth)
        \item Test each action individually
        \item Verify simulation correctness
    \end{itemize}
    
    \item \textbf{Phase 3: AI Implementation} (Week 3)
    \begin{itemize}
        \item Implement GreedyAI first (simplest)
        \item Create evaluation function
        \item Test AI vs random player
        \item Implement MinimaxAI
        \item Implement MCTS AI
    \end{itemize}
    
    \item \textbf{Phase 4: Visual Polish} (Week 4)
    \begin{itemize}
        \item Replace colored squares with tile assets
        \item Add animations (water, particles)
        \item Implement UI panel
        \item Add visual effects (glow, floating text)
        \item Polish and refine
    \end{itemize}
    
    \item \textbf{Phase 5: Testing and Tuning} (Week 5)
    \begin{itemize}
        \item Balance game parameters
        \item Test AI vs AI matches
        \item Collect performance data
        \item Create presentation materials
    \end{itemize}
\end{enumerate}

\subsection{Version Control and Collaboration}

Use Git for version control:

\begin{lstlisting}[language=bash]
# Initialize repository
git init
git add .
git commit -m "Initial project structure"

# Create branches for features
git checkout -b feature/ai-minimax
git checkout -b feature/visual-effects
\end{lstlisting}

\textbf{.gitignore file}:
\begin{lstlisting}
# Python
__pycache__/
*.pyc
*.pyo
.venv/
venv/

# IDE
.vscode/
.idea/
*.swp

# OS
.DS_Store
Thumbs.db

# Game data
saves/
logs/
\end{lstlisting}

% ================================================================
% CONCLUSION
% ================================================================
\section{Conclusion and Next Steps}

This document provides a comprehensive blueprint for creating "World Balance Architects," a strategic AI game suitable for laboratory demonstration and academic presentation.

\subsection{Key Strengths of This Design}

\begin{itemize}
    \item \textbf{AI-Friendly}: Turn-based, discrete actions perfect for AI comparison
    \item \textbf{Strategic Depth}: Simple rules create complex emergent gameplay
    \item \textbf{Visually Impressive}: Professional appearance without complex engines
    \item \textbf{Educational Value}: Demonstrates multiple AI algorithms clearly
    \item \textbf{Scope-Appropriate}: Achievable in 4--6 weeks for skilled developer
    \item \textbf{Extensible}: Easy to add features, modes, or AI variants
\end{itemize}

\subsection{Recommended Priorities}

\textbf{If time is limited, focus on}:
\begin{enumerate}
    \item Core simulation working correctly
    \item At least 2 AI levels (Greedy + Minimax)
    \item Clean, readable visuals (even if simple)
    \item Clear demonstration of AI comparison
\end{enumerate}

\textbf{If you have extra time, add}:
\begin{enumerate}
    \item MCTS AI implementation
    \item Advanced visual effects
    \item Multiple game modes (competitive, cooperative)
    \item Replay system for analyzing AI decisions
    \item Statistics and graphs showing AI performance
\end{enumerate}

\subsection{Presentation Tips for Laboratory}

When demonstrating to your instructor or class:

\begin{itemize}
    \item \textbf{Live Demo}: Show 2--3 turns of AI vs AI gameplay
    \item \textbf{Code Walkthrough}: Explain evaluation function and one AI algorithm
    \item \textbf{Performance Comparison}: Show graphs of AI win rates
    \item \textbf{Design Decisions}: Explain why you chose specific simulation rules
    \item \textbf{Challenges Overcome}: Discuss interesting bugs or optimization problems
\end{itemize}

\subsection{Resources for Further Learning}

\begin{itemize}
    \item \textbf{Pygame Documentation}: \url{https://www.pygame.org/docs/}
    \item \textbf{AI Algorithms}: "Artificial Intelligence: A Modern Approach" by Russell \& Norvig
    \item \textbf{Game AI}: \url{http://www.gameaipro.com/} (free online book)
    \item \textbf{Python Best Practices}: PEP 8 style guide
\end{itemize}

\vspace{1cm}

\begin{center}
\large
\textbf{Good luck with your project!}

\vspace{0.5cm}

This design balances ambition with feasibility, creating a game that will impress your instructors while remaining achievable within academic timeframes.
\end{center}

\end{document}
